\section{Conclusion}
\label{sec:conclusion}

This project asks whether India's overnight call-money rate can be
predicted at the weekly horizon from open public data, what drives
the prediction, and what the answer means for monetary policy
practice. The headline numbers are clean: 545 weekly observations
across 117 features assembled from ten public sources, two distinct
monetary regimes identified by K-Means at silhouette 0.4238, and a
walk-forward XGBoost forecast with RMSE 0.1019 (about ten basis
points) and 70.9\% directional accuracy across 389 out-of-sample test
weeks.

The substantive findings are stronger than the numbers alone suggest.
First, the structural break in early 2020 is real, persistent, and
visible in unsupervised clustering without any human labelling of
dates. Second, the policy corridor (Repo, Reverse Repo, MSF) plus the
autoregressive lags of WACMR account for essentially all the
predictive power; the twenty-eight equity and currency features added
through Yahoo Finance and custom technical indicators do not
materially help. Third, the model's tail errors are not driven by
predictable seasonal calendar pressure, which it learns well, but by
a small number of unsignalled inter-meeting policy actions, which by
construction it cannot.

The translation of these findings into action is straightforward. RBI
liquidity managers should monitor the WACMR-Repo spread as a
real-time barometer; treasury desks should size positions to the
model's regime confidence; equity and forex desks should stop
treating Nifty 50 and USD/INR technical indicators as inputs to
overnight-rate forecasts; and NDAP should publish weekly T-Bill
cut-off yields as a dedicated series.

Beyond the analytical findings, the project ships a fully deployed
interactive system. A Next.js dashboard on Vercel, a FastAPI backend
on Render with a Gemini 2.5 Flash function-calling agent, a
counterfactual simulator with a 26-second to 50-millisecond cache
optimisation, and a long-form blog with a live report. The system
runs on free-tier infrastructure, passes every Playwright audit page
with zero console errors, and serves the underlying data through a
typed REST API that any other researcher can build on.

Every figure in this report is reproducible from the code committed
to the repository. The pipeline stages are numbered, idempotent, and
deterministic. The SQLite database and the saved model artefacts are
committed alongside the code, so an evaluator can reproduce the
analysis without re-running the pipeline at all. The hope is that
this combination of substantive findings, working software, and full
reproducibility makes the project useful as both a research artefact
and a starting point for further work on Indian monetary-policy
analytics.
